\title{Large Language Models Can Automatically Engineer Features for Few-Shot Tabular Learning}

\begin{abstract}
Large Language Models (LLMs) have exhibited exceptional proficiency in solving complex and previously unencountered reasoning tasks, revealing substantial promise for applications in tabular learning—a domain critical across numerous practical settings. In this work, we introduce \textsf{FeatLLM}, an innovative in-context learning strategy that leverages LLMs for automatic feature engineering, thereby crafting input datasets highly tailored for tabular predictive tasks. These engineered features are then fed into a basic downstream machine learning algorithm, such as linear regression, leading to robust few-shot learning outcomes. Notably, during inference, \textsf{FeatLLM} relies exclusively on this elementary predictive model and the generated features, forgoing the need for additional LLM queries for every test sample, unlike most existing LLM-based methods. This approach requires only API-level LLM access, surmounts limitations imposed by prompt length, and reduces computational demands at inference. Through extensive evaluation on a broad suite of tabular benchmarks, \textsf{FeatLLM} consistently delivers superior rule-based feature generation and yields considerable (averaging a 10\% margin) gains over comparable systems such as TabLLM and STUNT.
\end{abstract}

\section{Introduction}
Large Language Models (LLMs) have recently demonstrated remarkable capability to generate responses for tasks not seen during training, doing so without necessitating tailored fine-tuning for each new task~\cite{creswell2022selection,imani2023mathprompter,taylor2022galactica}. Their vast exposure to large-scale text corpora enables LLMs to encode extensive background knowledge, which often compensates for domain-specific expertise~\cite{Genomes2021,singhal2023large,wilcox2003role}, thus proving invaluable in automating a spectrum of tasks from dialogue understanding~\cite{finch2023leveraging} to software debugging~\cite{fan2023automated}. By supplementing prompts with descriptions of the target task and providing a handful of illustrative examples, LLMs are able to infer contextual cues and selectively focus on salient attributes and instances~\cite{brown2020language,zhao2023survey}. This underlines their capacity not only for comprehension, but also for engaging in advanced data analysis and expert-level feature construction.

The application of LLMs’ foundational knowledge and sophisticated reasoning is increasingly being adapted to tabular learning settings. For example, leveraging external knowledge such as age-related disease risk can substantially improve performance in medical prediction scenarios. Prevailing methodologies tend to articulate task goals and feature meanings in natural language, represent tabular data in serialized forms, and then submit this composite input to an LLM~\cite{dinh2022lift,wang2023anypredict}. This is commonly paired either with an in-context learning framework~\cite{nam2023semi} or with efficient adaptation via fine-tuning~\cite{dinh2022lift,hegselmann2023tabllm}. Such techniques have shown that the explicit integration of LLMs’ knowledge routinely surpasses conventional tabular modeling baselines, particularly in data-scarce conditions~\cite{hegselmann2023tabllm}.

Existing approaches that utilize LLMs for tabular data modeling are subject to several drawbacks. Achieving end-to-end predictions generally necessitates running an LLM inference on every single data point, leading to significant computational overhead. High-precision results are typically attainable only when the LLM can be fine-tuned, an option that is frequently unavailable for many state-of-the-art models that restrict users to inference-only APIs and do not provide public access to training resources~\cite{anil2023palm,openai2023gpt4}. Another considerable issue is the inadequate performance of current strategies with extensive prompt lengths~\cite{liu2023lost}; as the number of variables or features increases and their textual encoding exceeds manageable limits, both feasibility and model effectiveness sharply decline. Collectively, these issues present major obstacles for deploying LLM-driven tabular models in practical settings.

To address these challenges, our method reframes the role of LLMs away from direct end-to-end prediction and instead emphasizes their use in the feature engineering phase. Concretely, our in-context learning framework, \textsf{FeatLLM}, is designed to infer the underlying ``criteria'' or decision mechanisms that inform the model’s predictions. For example, in a medical diagnosis task, the LLM is prompted to deduce and articulate explicit rules indicating which feature combinations are indicative of the presence of a disease. These decision criteria are then leveraged to construct synthetic features that supplant the original ones during model training. As a result, the core function of the LLM is shifted to automated feature creation, leading to a simplified downstream modeling process; a lightweight machine learning model can deliver predictions based on the newly engineered features, substantially reducing inference costs. By capitalizing on in-context learning, features are extracted without any LLM fine-tuning—simply through strategically crafted prompt demonstrations. Moreover, methodologies inspired by ensemble learning, such as feature bagging~\cite{breiman2001random}, are utilized to control and minimize prompt length, mitigating the adverse effects associated with very large input texts.

\textsf{FeatLLM} organizes the feature engineering prompt into two primary tasks: (1) comprehending the underlying problem and inferring associations between features and target variables; and (2) leveraging this understanding along with a limited number of examples to formulate discriminative rules for class separation. This systematic reasoning framework enables LLMs to disregard irrelevant features during the process of rule generation. The extracted rules are subsequently enacted on the dataset—where the data are interpreted as programmatic representations—resulting in new binary features that indicate whether individual instances satisfy each rule. Subsequently, a simple model is trained using these features to predict class membership probabilities. Robustness is further bolstered by an ensemble methodology, involving repeated feature generation combined with feature bagging. By calibrating the feature and sample selection during bagging, the issue of overly large input prompts is alleviated. Since \textsf{FeatLLM} operates without training overhead and incurs minimal inference costs, it enables large language model applications to diverse real-world contexts. 

Our assessment of \textsf{FeatLLM} spans 13 tabular datasets under low-data scenarios, demonstrating its consistent and impressive effectiveness. The results confirm that LLMs can harness both prior and data-driven insights, generating features of substantial quality. Across various experimental conditions, our approach consistently surpasses established few-shot learning baselines. The associated code is made available via an anonymized GitHub repository at~\texttt{https://github.com/Sungwon-Han/FeatLLM}.

\section{Related Work}

\subsection{Few-Shot Learning with Tabular Data} Few-shot learning has emerged as a powerful technique for deriving transferable representations from datasets, while substantially reducing the need for extensive labeling~\cite{chen2018closer}. While initial research prioritized vision tasks, few-shot methods have exhibited robust generalization capacities beyond image data, expanding into domains such as text, audio, and, increasingly, tabular data~\cite{majumder2022few,nam2023stunt,schick2021exploiting}. However, extracting reliable patterns from a limited labeled dataset can result in the identification of spurious patterns or artifacts~\cite{bartlett2021deep}. This challenge has motivated recent works to incorporate supplementary information and embed domain-specific prior knowledge, thus promoting beneficial inductive biases during training. Strategies to address data scarcity include leveraging vast repositories of unlabeled tabular datasets and employing strategic data augmentation within semi-supervised frameworks~\cite{ucar2021subtab,yoon2020vime}, or employing task-independent unsupervised pre-training methods for improved model initialization~\cite{somepalli2021saint}. Typical unsupervised learning paradigms might utilize strategies such as data masking or corruption, optimizing a reconstruction objective~\cite{arik2021tabnet,majmundar2022met}, or applying data augmentations suited for contrastive learning~\cite{bahri2022scarf,chen2020simple}. Nonetheless, the diverse and non-uniform properties of tabular data limit the efficacy and universality of such augmentations, and these approaches have demonstrated restricted benefits in few-shot scenarios~\cite{nam2023stunt}. 

To circumvent these limitations, contemporary research advocates for pre-training models on a wide assortment of tabular datasets, many of which are not closely aligned to the target tasks, and subsequently transferring the learned representations~\cite{zhang2023generative,levin2022transfer,zhu2023xtab}. One notable example, TransTab~\cite{wang2022transtab}, adopts a transformer-based architecture that captures semantic relationships among columns, inferred both from column names and their contents. Knowledge accumulated from this diverse corpus of tables has enabled models to perform zero-shot or few-shot generalization on novel tabular problems. Similarly, TabPFN~\cite{hollmann2022tabpfn} generates mixed synthetic distributions for pre-training a Bayesian neural network. For inference, both the training and test data from a target domain are provided simultaneously, negating the need for further model updates. The assimilated prior across diverse data sources allows these approaches to outperform conventional tree-based models, especially in data-scarce scenarios.

\subsection{Language-Interfaced Tabular Learning} A further avenue for encoding external knowledge into tabular models harnesses the capabilities of large language models (LLMs)~\cite{anil2023palm,openai2023gpt4}. Benefiting from training on substantial and varied textual corpora, LLMs exhibit an aptitude for extrapolating to unseen tasks and transferring knowledge across distinct contexts~\cite{brown2020language}. Recent investigations have concentrated on channeling this latent linguistic knowledge toward tabular learning tasks. The prevalent methodology involves transforming tabular data into textual sequences for integration into LLM input prompts~\cite{dinh2022lift,hegselmann2023tabllm,wang2023anypredict}. By embedding table contents, descriptive metadata, and task instructions within prompts, LLMs can model associations between task requirements and feature spaces to enable inference. These methods may further enhance performance by adopting parameter-efficient tuning approaches such as LoRA~\cite{hu2021lora} or IA3~\cite{liu2022few}, or by utilizing in-context learning paradigms where labeled examples are included directly within prompts, allowing models to generalize without fine-tuning~\cite{dinh2022lift,hegselmann2023tabllm,nam2023semi,slack2023tablet}.

While previous methods have shown success in enhancing few-shot tabular learning, their reliance on conducting inference solely through the LLM introduces obstacles for real-world industrial implementation. This work seeks to break from the traditional paradigm of using LLMs as end-to-end black-box predictors for individual samples. Instead, the focus is on leveraging the LLM to infer explicit class-specific rules or conditions. \textsf{FeatLLM} takes the rules produced by the LLM, translates them into programmatic code to construct new feature representations, and thereby permits prediction without continually routing data through the LLM. In this approach, in-context learning is harnessed to extract meaningful rules by drawing on pre-existing knowledge as well as information from the few-shot exemplars.

\section{Methods}
\textsf{FeatLLM} operates by identifying ``rules"—that is, the distinctive conditions that delineate each possible class—using the limited set of input samples, as illustrated in Figure~\ref{fig:main_model}. After these rules are articulated via the LLM, they are interpreted as executable code, subsequently yielding novel binary features for each instance to signal if a rule has been satisfied. These binary features are then combined with a set of non-negative, learnable weights through a dot product to estimate the probability that the sample belongs to a given class. Through model optimization, one can uncover which rules (features) most influence the predicted class probabilities (see Section~\ref{Sec:training}). The overall procedure leverages multiple iterations utilizing bagging, and their outputs are merged using ensemble techniques. Details concerning problem formalization and the specific steps of each phase follow in the next sections.

\vspace*{-0.12in}
\paragraph{Problem formulation.} Consider a tabular dataset comprising $N$ labeled instances, denoted as $\mathcal{D} = \{(\mathbf{x}^i, \mathbf{y}^i)\}_{i=1}^{N}$. Each sample $\mathbf{x}^i$ is represented as a $d$-dimensional feature vector, and the feature names are given as a set $F = \{f_j\}_{j=1}^{d}$, each accompanied by their respective natural language definitions. Within a classification framework, every label $\mathbf{y}^i$ takes a value from a predefined collection of classes. For $k$-shot learning scenarios, the model is trained using only $k$ ($<N$) labeled instances selected randomly.

\subsection{Prompt Design for Extracting Rules}
\label{Sec:prompt_design}
To ensure that the LLM derives rules via a reasoning trajectory similar to that of expert practitioners addressing tabular learning challenges, we carefully structure the prompt to encourage such cognitive processes. Our prompt strategy includes three fundamental elements (see Fig.~\ref{fig:prompt_for_rule} and Appendix~\ref{sec:example_rule}):

\vspace*{-0.12in}
\paragraph{Basic information description.} This section communicates all necessary background for approaching the given task (highlighted in orange in Fig.~\ref{fig:prompt_for_rule}). It outlines the problem statement, presents the label definitions, offers feature details and, where appropriate, their explanations, and introduces illustrative examples. The primary task is posed as a direct question (e.g., ``Does the data on this patient's myocardial infarction complications suggest chronic heart failure? Yes or no?"; refer to Appendix~\ref{sec:task_desc} for additional instances). Each feature’s value type is specified, and, when available, the definition is appended; for categorical features, sample values for each category are given. For demonstration, a selected subset of training samples is linearized as textual sequences paired with their corresponding labels, adhering to the textual formatting conventions adopted in previous works~\cite{dinh2022lift,hegselmann2023tabllm}:

\begin{align}
&\text{Serialize}(\mathbf{x}^i,\mathbf{y}^i, F) = \nonumber \\ 
&``\ f_1\ \text{is}\ \mathbf{x}^i_1.\ ... \ f_d\  \text{is}\  \mathbf{x}^i_d.\ \ \text{Answer:}\  \mathbf{y}^i\ ”
\end{align}

\vspace*{-0.12in}
\paragraph{Reasoning instruction.} Our objective is to strengthen the reasoning capabilities of the LLM by supplying it with reasoning-focused instructions (refer to the blue text in Fig.~\ref{fig:prompt_for_rule}). The reasoning instruction commences with an opening sentence inspired by the chain-of-thought methodology~\cite{wei2022chain}. Subsequently, we partition the LLM’s inference process into two distinct phases. Initially, the LLM is prompted to deduce the causal relationships or inherent tendencies between features and the overarching task description, deliberately excluding any reference to example demonstrations. This step depends exclusively on broad-based knowledge or commonsense reasoning, enabling the LLM to capitalize on its internalized knowledge relevant to the task. In the subsequent phase, the LLM utilizes both the example demonstrations and its preliminary insights to derive classification rules specific to each class. Implementing this two-stage reasoning minimizes the risk of the LLM focusing on coincidental patterns in irrelevant fields and assists in prioritizing key features. In order to avoid generating too few rules (which could result in underfitting) or too many (leading to unwieldy prompts), we standardize the process by setting the number of rules to extract per class at a fixed quantity (i.e., 10)\footnote{Further analysis on rule number selection can be found in Section~\ref{sec:analysis}.}. Furthermore, to prevent type inconsistencies in rule construction, the LLM is directed to consult the feature type details provided within the general information segment. 

\vspace*{-0.12in}
\paragraph{Response instruction.} To make the output rules easier to parse and apply, we provide the LLM with detailed guidance on formatting its responses (refer to the yellow text in Fig.~\ref{fig:prompt_for_rule}). The prompt specifies the expected format for each answer class (i.e., how each response should be structured) along with the precise formatting to be followed for individual rules (i.e., format for the rule). Such instructions empower the LLM to choose formats that align with different feature types. If no explicit guidance is present, the LLM may autonomously integrate multiple rules using logical connectives such as AND or OR. An illustrative result is available in Appendix~\ref{sec:outcome-rule}.

\subsection{Inferring Class Likelihood via Rules}
\label{Sec:training}

\paragraph{Parsing rules for feature generation.} The rules obtained from the earlier step are used to synthesize new binary features. For each target class, these features encode whether a particular sample meets the criteria specified by the rules for that class (i.e., a binary value of '0' or '1'). Such features can aid in assessing a sample's likelihood of membership in any given class. Nonetheless, since the LLM produces rules in natural language, a specific parsing strategy must be developed to automate the generation of these features. Even with enforced constraints on the rule generation format for improved parsability, the LLM's outputs frequently exhibit minor syntactic variations or noise~\cite{kovavc2023large}. As an example, expressions such as ``$>$” and ``$<$” can be reformulated as ``greater than" or ``less than" within the LLM output, further complicating rule parsing.

To overcome the challenge of extracting consistent features from noisy or variably formatted text, rather than developing intricate custom parsers, we exploit the strengths of the LLM itself. The LLM is adept at interpreting the intended semantics in text despite variations in surface structure, and can generate simple executable code based on provided instructions, as its training includes code-related corpora. To harness this capability, prompts incorporating the function signature, as well as descriptions of input and output formats and the applicable rules, are constructed and submitted to the LLM. Representative prompts and resulting code, as detailed in Figures~\ref{fig:prompt_for_parsing} and ~\ref{sec:example_parsing}-\ref{sec:example_parse_outcome} of the Appendix, illustrate this approach. The generated code is then applied for data transformation through Python’s \texttt{exec()} functionality.

\vspace*{-0.12in}
\paragraph{Inferring class likelihood.} After extracting binary rule-based features for every class, an initial, straightforward way to estimate a sample's propensity for belonging to a specific class is to tally the number of satisfied rules per class (i.e., summing binary indicators within each class). Nevertheless, these rules and their resultant features can differ in their predictive relevance, making it important to empirically determine their relative weights using labeled training data. This is approached by fitting a standard machine learning (ML) model, focused on the binary feature vector associated with each class. For illustration, consider employing a bias-free linear model. Let $\mathbf{z}_k^i \in \{0, 1\}^R$ represent the binary features for class $k$ derived from input sample $\mathbf{x}^i$, where $R$ is the count of rules for each class and $c$ refers to the total number of classes. The estimated probability vector for each class, $\mathbf{p}^i$, is given by:

\begin{align}
\text{logit}_k^i &= \max(\mathbf{w}_k, 0) \cdot \mathbf{z}_k^i, \nonumber \\
\mathbf{p}^i &= \text{Softmax}({[\text{logit}_{1}^i, ..., \text{logit}_{c}^i]}). \label{eq:infer_prob}
\end{align}

In this context, $\mathbf{w}_k$ denotes the trainable weights within the linear model, which quantify the significance of each corresponding feature. By constraining these weights to reside in the positive domain, it is guaranteed that the features generated by the LLM cannot negatively influence a class’s probability during the training phase. Following this, the output logit vector is transformed into a set of class probabilities by applying the softmax operation. The weight parameters $\mathbf{w}_k$ are learned by minimizing the cross-entropy loss using a limited set of training instances in a few-shot regime.

\vspace*{-0.12in}
\paragraph{Ensembling with bagging.} The complete procedure is carried out multiple times to generate an ensemble of inference models, whose predictions are aggregated to yield the final output. For $T$ independent iterations, the resulting trained weights, noted as $\{\mathbf{w}[t]\}_{t=1}^T$, are used to obtain class probability predictions. The final class probability $\mathbf{p}^i$ is obtained by averaging the probabilities produced from each $\mathbf{w}[t]$ according to Eq.~\ref{eq:infer_prob}.

To introduce sufficient diversity for effective ensemble learning, a high temperature setting is used during LLM rule generation. The order of few-shot exemplars is also randomized, informed by observations that the demonstration sequence can influence LLM outputs~\cite{lu2022fantastically}\footnote{Refer to Appendix~\ref{sec:diversity} for an analysis of rule diversity.}. Bagging is further employed to select a random subset of features or samples in each iteration, which becomes especially beneficial when the training set or feature set is large. This ensemble methodology provides two significant benefits: firstly, any errors produced by faulty rules in individual runs can be mitigated by accurate rules from other runs, collectively boosting robustness through a majority-consensus mechanism that resonates with the concept of LLM self-consistency~\cite{wang2022self}; secondly, this approach counters constraints posed by the limited context window of LLMs—bagging helps manage prompt length when input data exceeds model limits, thereby maintaining performance stability. Although any single set of rules may cover only part of the feature space, aggregating numerous weak learners derived from different runs compensates for incomplete coverage and enhances model reliability.

\section{Experiments}
We assess the effectiveness of \textsf{FeatLLM} on a diverse collection of tabular datasets and undertake multiple analyses to gain insights into our framework's behavior. Owing to limitations on space, further experimental results—including studies on different LLMs, qualitative evaluations, and supplementary analyses—are presented in the Appendix.

\subsection{Performance Evaluation}
\paragraph{Datasets.}
Our experimental evaluation spans 13 datasets designed for binary or multi-class classification. The selected datasets include: (1) Adult~\cite{asuncion2007uci}, which aims to predict whether an individual’s income surpasses \$50,000 annually; (2) Bank~\cite{moro2014data}, focusing on term deposit subscription prediction; (3) Blood~\cite{yeh2009knowledge}, which predicts the likelihood of donor return; (4) Car~\cite{kadra2021well}, which classifies car quality; (5) Communities~\cite{misc_communities_and_crime_183}, targeting crime rate estimation within a community; (6) Credit-g~\cite{kadra2021well}, assessing whether an individual is likely a favorable or unfavorable credit risk; (7) Diabetes\footnote{\texttt{kaggle.com/datasets/uciml/pima-indians-diabetes-database}}, predicting diabetes status for an individual; (8) Heart\footnote{\texttt{kaggle.com/datasets/fedesoriano/heart-failure-prediction}}, which forecasts coronary artery disease occurrence; (9) Myocardial~\cite{misc_myocardial_infarction_complications_579}, focusing on chronic heart failure prediction.

Additionally, we have incorporated newer and synthetic datasets to examine generalization beyond commonly used benchmarks and reduce potential overlap with LLM pre-training data: (10) Cultivars, concerned with estimating the grain yield class of a particular soybean cultivar\footnote{\texttt{archive.ics.uci.edu/dataset/913}}; (11) NHANES, where the goal is to classify an individual's age group given tabular features\footnote{\texttt{archive.ics.uci.edu/dataset/887}}; (12) Sequence-type, which determines the category of a sequence from arithmetic, geometric, Fibonacci, or Collatz; (13) Solution-mix, which predicts whether a mixture of four solutions exceeds a concentration threshold of 0.5. The first two datasets were contributed to the UCI Machine Learning Repository post-September 2023, ensuring exclusion from the pre-training corpus of the LLM API (gpt-3.5-turbo-0613) utilized. The latter two represent synthetic datasets crafted to avoid possible memorization by the LLM, thus providing a robust benchmark for our approach.

These datasets span a spectrum in terms of size and intrinsic difficulty, as reported in Table~\ref{Tab:basic-desc}. Each includes well-defined, descriptive feature names and attribute explanations. Datasets such as ‘Communities’ and ‘Myocardial’ present scalability challenges for LLMs, given their high-dimensional input with over 100 features and limitations stemming from restricted context window size.

\vspace*{-0.12in}
\paragraph{Baselines.}
We benchmark our method against nine alternative approaches. The first trio represents traditional algorithms for tabular learning: (1) Logistic Regression (LogReg), (2) XGBoost~\cite{chen2016xgboost}, and (3) RandomForest~\cite{ho1995random}. The subsequent two methods, (4) SCARF~\cite{bahri2022scarf} and (5) STUNT~\cite{nam2023stunt}, exploit unlabeled data; it should be noted, however, that access to extensive unlabeled data is often impractical in real settings. The sixth competitor, (6) TabPFN~\cite{hollmann2022tabpfn}, leverages synthetic datasets with varied distributions during pretraining. The final group employs LLM-driven strategies: (7) In-context learning (In-context)~\cite{wei2022emergent}, (8) TABLET~\cite{slack2023tablet}, and (9) TabLLM~\cite{hegselmann2023tabllm}. In-context learning directly introduces a few labeled training instances as prompts to the LLM without parameter adaptation. TABLET supplements these prompts with additional information—via rule-based hints and exemplar prototypes sourced from external classifiers—to refine model inference. TabLLM utilizes parameter-efficient fine-tuning strategies such as IA3~\cite{liu2022few}, adapting LLMs through training on limited labeled samples.

\vspace*{-0.12in}
\paragraph{Implementation details.}
\textsf{FeatLLM} adopts GPT-3.5 as its primary LLM architecture, but it is built to function independently of any specific LLM, as evidenced by experiments with alternative backbones provided in Appendix~\ref{sec:backbones}. The LLM inference is conducted with a temperature of 0.5, and the top-p parameter retains its default API value of 1. We configure the ensemble count to 20 and extract 10 rules per instance. Further analysis of the hyper-parameter settings and their effects can be found in Figure~\ref{fig:hyperparameter2} as well as Appendix~\ref{sec:hyper-parameter}.

For the linear classifier, we utilize the Adam optimization algorithm with a learning rate fixed at 0.01, proceeding for 200 epochs. Model selection is performed via $k$-fold cross-validation. Baseline reproduction makes use of GPT-3.5 for in-context learning scenarios, while TabLLM is evaluated with T0~\cite{sanh2022multitask} as stipulated in its original protocol. Importantly, TabLLM mandates that only LLMs with open checkpoints may be used, thus excluding GPT-3.5 from these experiments. Standard machine learning algorithms such as LogReg, XGBoost, and RandomForest have their hyper-parameters tuned by executing a Grid-search in combination with $k$-fold cross-validation. The parameter $k$ takes a value of either 2 or 4, with the constraint that the training partition in each fold contains at least one sample from each class. Additional implementation specifics can be located in Appendix~\ref{sec:baselines}.

\vspace*{-0.12in}
\paragraph{Results.}
Tables~\ref{tab:main1} and \ref{tab:main2} display comparative results for a diverse collection of datasets. Each experiment was conducted three times, varying the training and testing data selections, with the means and standard deviations for AUC scores reported. \textsf{FeatLLM} reliably achieves either the best or second-best outcomes among all baseline methods. Remarkably, with just 4-shot inputs, our method approaches the effectiveness of traditional tabular approaches that typically necessitate 32 shots, as depicted in Fig.~\ref{fig:compares}. While STUNT and TabPFN yielded notable gains on a select group of datasets, their aggregate improvement over standard machine learning approaches remained limited. Additionally, LLM-based methods including in-context learning, TABLET, and TabLLM encountered difficulties and were unable to process tabular datasets with extensive feature sets. Our approach, which incorporates both bagging and ensemble mechanisms, maintained competitive accuracy on high-dimensional datasets. This highlights its robustness. It should also be considered that while our bagging and ensembling paradigm may be transferable to other LLM-centric methodologies, implementing this would result in a twofold increase in per-sample inference workload, thereby markedly increasing computational demands to the point of being infeasible in practice.

\subsection{Analysis \& Discussion}
\label{sec:analysis}
\paragraph{Ablation study.} To evaluate the role of key components in overall model effectiveness, we performed a series of ablation studies, examining: (1) \textsf{-Tuning}: excluding the weight adjustment stage in the linear model, opting instead to aggregate the binary features for logit determination; (2) \textsf{-Ensemble}: removing the ensemble procedure; (3) \textsf{-Description}: excluding feature-level descriptions; and (4) \textsf{-Reasoning}: not including the initial reasoning step (Step-1) when formulating rules.

Table~\ref{tab:ablation} provides the average changes in AUC following each ablation, aggregated across all datasets. In every scenario, omitting these elements leads to worsened performance. Of particular interest, the advantage conferred by weight tuning becomes more pronounced as the number of shots rises; when abundant examples are available, the precise evaluation of rule importance is more achievable, thus amplifying the tuning effect. Conversely, the presence of feature descriptions and reasoning procedures becomes especially advantageous with fewer shots, emphasizing that leveraging the prior knowledge embedded in LLMs is essential in low-data regimes. Finally, our ensemble strategy consistently contributes to performance gains in all configurations.

\vspace*{-0.12in}
\paragraph{Training \& inference time comparison.} We evaluate and contrast the time requirements for training and inference across various methods. Our experiments utilize the Adult dataset with a 16-shot training set. Unless otherwise specified, a single A100 GPU is employed, while TabLLM necessitates four A100 GPUs to facilitate model parallelism. Table~\ref{tab:cost} reports detailed runtime figures. A key observation is that \textsf{FeatLLM} achieves inference speeds that are relatively low and on par with traditional approaches. Conversely, alternative LLM-based techniques—including In-context learning, TABLET, and TabLLM—incur significantly longer inference durations. Regarding the training phase, \textsf{FeatLLM} aligns closely with other LLM-based models in terms of time efficiency: it issues only 30 API calls, which can be executed in parallel and does not cause notable budgetary overhead.

\vspace*{-0.12in}
\paragraph{Quality of features generated by \textsf{FeatLLM}.}
We perform a basic assessment to understand the informativeness of the rule-based features produced by \textsf{FeatLLM} when addressing real-world tasks. Concretely, we compute the mean AUROC between the features generated and the target labels and determine the proportion of features deemed useful, specifically those with an AUROC exceeding 0.5, across different datasets. 
The outcomes, summarized in Table~\ref{tab:quality}, demonstrate that most of the constructed rule features are indeed pertinent and offer meaningful guidance for solving the target problem (see the “Before tuning” column). Additionally, when the linear model undergoes tuning with actual data, those rules that contribute little—specifically, those with assigned weights below a threshold—are automatically discarded (refer to the “After tuning” column). Here, the threshold is set at 0.1, established based on the distribution of all weight values.

\vspace*{-0.12in}
\paragraph{Effect of adding spurious correlations.} To evaluate the capabilities of different approaches in excluding irrelevant spurious correlations under limited data settings, we performed a systematic study. Specifically, we worked with the Heart dataset, augmenting it by randomly appending feature columns from the Adult dataset, which hold no relevance to the Heart task. We then tracked model performance as we sequentially increased the quantity of these extraneous features. To ensure fairness, we standardized the number of labeled examples at 8 and did not involve any unlabeled data, excluding algorithms such as SCARF and STUNT from the set of baselines.

Figure~\ref{fig:spurious} provides a graphical summary of how additional spurious correlations impact model accuracy. In these tests, \textsf{FeatLLM} consistently experienced the smallest decline in performance, outperforming the other baselines. This relative resilience can be attributed to our framework’s reliance on prior knowledge for establishing task-feature connections during rule derivation, which discourages creating rules around unrelated attributes and thus supports robust predictions. Empirical analysis confirms this, as rules derived from noisy columns comprised only 1-2\% of the total set. Nevertheless, \textsf{FeatLLM} is not immune to mistakenly modeling spurious correlations when additional features from similar domains are included at random; we provide further examination of such scenarios in Appendix~\ref{sec:spurious-appendix}.

\vspace*{-0.12in}
\paragraph{Hyper-parameter analysis}
\textsf{FeatLLM} operates with two principal hyper-parameters: the ensemble size, and the number of rules per class produced in each LLM invocation. We systematically explored the influence of these hyper-parameters on the overall model performance. Our findings indicate that performance generally rises as the ensemble count is increased; however, after reaching approximately 20 ensembles, the benefits plateau (refer to Fig.~\ref{fig:appendix-hyperparameter1} in the Appendix). Regarding the count of rules per query, there was no statistically meaningful effect detected, yet, considering both prompt length and predictive accuracy, 10 emerged as an optimal setting (see Fig.~\ref{fig:hyperparameter2}). We hypothesize that introducing additional rules escalates the input's dimensionality, which injects more information but also raises the risk of overfitting in low-shot settings.

\section{Conclusion}
We introduce \textsf{FeatLLM}, which sets a new benchmark in leveraging LLMs for few-shot tabular learning. Departing from conventional methods that employ LLMs for direct sample-wise predictions, \textsf{FeatLLM} generates predictive heuristics akin to those employed by feature engineers. This rule-induction mechanism, together with bagged ensembles, yields minimal inference costs, functions with simple API interactions (bypassing full retraining), and is agnostic to data dimensionality constraints. Empirically, \textsf{FeatLLM} demonstrates robust predictive accuracy, offering a compelling, data-efficient LLM approach for practical tabular datasets.

While \textsf{FeatLLM} is tailored for scenarios with limited labeled data, our prospective work seeks to broaden its scope: crafting novel features for large-sample datasets, experimenting with alternative feature types beyond rules, and enhancing interpretability—thus extending its applicability to a wider spectrum of tasks.

\section{Broader Impact}
The proposed approach, \textsf{FeatLLM}, is motivated by leveraging LLM-derived prior knowledge to surpass what conventional supervised models can achieve in tabular settings. By maximizing learning efficiency, our work plays a crucial role in addressing overfitting risks associated with scant training data through the integration of prior knowledge. \textsf{FeatLLM} is particularly advantageous for practitioners facing limited labeled data—such as in finance or healthcare, where annotation is expensive or impractical, and LLMs' domain-relevant pretraining can provide substantial value. Moving forward, a vital consideration for wider deployment involves systematically examining the repercussions of societal biases and misinformation embedded in LLM pretraining, as such priors may disproportionately steer predictions, occasionally superseding available labeled data.

\end{document}